\begin{table}[htbp]\centering
\def\sym#1{\ifmmode^{#1}\else\(^{#1}\)\fi}
\caption{稳健性检验\label{tab:robustness}}
\begin{tabular}{l*{6}{c}}
\toprule
                    &\multicolumn{1}{c}{(1)}&\multicolumn{1}{c}{(2)}&\multicolumn{1}{c}{(3)}&\multicolumn{1}{c}{(4)}&\multicolumn{1}{c}{(5)}&\multicolumn{1}{c}{(6)}\\
                    &\multicolumn{1}{c}{基准}&\multicolumn{1}{c}{替换D}&\multicolumn{1}{c}{剔除直辖市}&\multicolumn{1}{c}{对数变换}&\multicolumn{1}{c}{剔除COVID}&\multicolumn{1}{c}{5\%缩尾}\\
\midrule
digital\_economy\_index&       0.277\sym{***}&                     &       0.142\sym{*}  &       0.252\sym{***}&       0.473\sym{***}&       0.242\sym{***}\\
                    &     (0.085)         &                     &     (0.070)         &     (0.077)         &     (0.136)         &     (0.063)         \\
\addlinespace
digital\_finance\_index&                     &       0.055         &                     &                     &                     &                     \\
                    &                     &     (0.037)         &                     &                     &                     &                     \\
\addlinespace
ln(人均GDP)         &                     &                     &                     &      -0.093         &                     &                     \\
                    &                     &                     &                     &     (0.064)         &                     &                     \\
\midrule
Observations        &         403         &         403         &         364         &         403         &         372         &         403         \\
\bottomrule
\multicolumn{7}{l}{\footnotesize 注：所有模型均包含控制变量、省份固定效应和年份固定效应，省份聚类标准误。}\\
\multicolumn{7}{l}{\footnotesize R1基准模型，R2替换核心解释变量为数字金融指数，R3剔除北京/天津/上海三个直辖市。}\\
\multicolumn{7}{l}{\footnotesize R4对数变换人均GDP和教育支出，R5剔除2020年，R6使用5\%缩尾替代1\%。}\\
\multicolumn{7}{l}{\footnotesize * p$<$0.10, ** p$<$0.05, *** p$<$0.01}\\
\end{tabular}
\end{table}
